\documentclass[11pt]{article}

\usepackage[utf8]{inputenc}
\usepackage[T1]{fontenc}
\usepackage[margin=1in]{geometry}
\usepackage{amsmath,amssymb}
\usepackage{caption}
% ponytail: this minimal TeX install lacks biblatex/biber and even
% hyperref's etoolbox dependency, so in-text citations below are
% hand-typed plain text rather than \textcite/\parencite calls, and
% cross-references use plain \ref (no clickable links). references.bib
% is kept as the structured source of truth for the reference list even
% though nothing in this .tex file parses it directly.

% ponytail: title page is a placeholder — replace with the course's official
% template before submission (user asked to defer this to a later pass).
\title{Does Remote Work Narrow the Motherhood Penalty?\\
Evidence from the Israeli Labor Market After COVID-19}
\author{[Student Name] \\{} [Course Name, Instructor] \\{} [Student ID]}
\date{\today}

\begin{document}

\maketitle

\begin{abstract}
\noindent The motherhood penalty is the well-documented drop in women's earnings and employment after childbirth. It persists for years and has no counterpart for fathers. Workplace inflexibility is a leading explanation for why it does not close. The COVID-19 pandemic's rapid shift toward remote work changed the constraints under which working mothers supply labor, offering a natural test of whether WFH availability narrows the penalty. This paper asks whether occupation-level exposure to work-from-home (WFH) opportunities differentially affected mothers' employment and hours in Israel, a labor market combining high female labor-force attachment with the OECD's highest fertility rate. The analysis uses Israeli Central Bureau of Statistics Labor Force Survey microdata (2017--2019, 2021--2023) and a triple-difference design built on a pre-determined, occupation-based WFH-exposure measure. I find no detectable narrowing of the penalty's extensive-margin (employment) component in the pooled post-COVID period. Mothers' hours worked, however, rise robustly and meaningfully conditional on employment (+0.83 hours/week, $p<0.01$). The triple-difference linking either margin specifically to WFH exposure is statistically underpowered rather than a clean zero. A systematic search across five extensive-margin heterogeneity dimensions turns up exactly one robust effect: employment among mothers of children aged 0--4 fell relative to other mothers after 2021. This pattern survives multiple robustness checks but cannot itself be linked to WFH exposure. The overall pattern is consistent with the motherhood penalty's post-pandemic adjustment concentrating on the cheaper, intensive margin, while Israel's furlough (\textit{Halat}) institution keeps the formal employment relationship largely intact.
\end{abstract}

\section{Introduction}

This paper examines whether the pandemic-driven reallocation of work toward the home altered the labor-market position of mothers relative to other workers. The question carries both scientific and policy significance. Scientifically, the sudden, large, and plausibly exogenous shift in the location of work constitutes a rare natural experiment in workplace flexibility. This is the very margin that Goldin (2014) identifies as the last major driver of the remaining gender pay gap in developed labor markets. Temporal and locational flexibility is a good mothers may pay for, in the form of lower wages or slower promotion. A shock that lowers the cost of supplying flexibility should, in principle, be detectable in mothers' labor-supply choices. Unlike the acute, transitory disruption of 2020, work-from-home (WFH) has proven to be a persistent feature of the post-pandemic labor market rather than a temporary accommodation (Barrero, Bloom, and Davis, 2021). This renders the question policy-relevant rather than merely historical.

The policy stakes are particularly high in Israel, which combines a highly productive, internationally integrated labor market with the OECD's highest total fertility rate. A labor market that induces mothers of young children to reduce employment or hours imposes a substantial aggregate cost. This cost is especially large in a country where a large and growing share of the working-age population is, at any given time, a parent of young children. Conversely, flexible work arrangements might measurably narrow the gap between mothers and other workers. If so, the finding bears directly on the design of remote-work policy, parental-leave rules, and public discourse on gender equality in the labor market. The remainder of this paper develops an empirical design that isolates the causal role of WFH \emph{exposure} -- as distinct from realized, endogenously chosen WFH status -- in shaping mothers' employment and hours after 2021.

\section{Background and Literature Review}

Three literatures bear on this question, and the contribution of this paper is best understood at their intersection.

\paragraph{The motherhood penalty and workplace flexibility.} A well-established literature documents a persistent earnings and employment gap that opens after the birth of a first child and does not close. Using Danish administrative data, Kleven, Landais, and Søgaard (2019) show that women's earnings fall by roughly 20\% relative to a pre-birth trend and never fully recover, while men's earnings are essentially unaffected. This is the canonical ``child penalty'' event-study design. Kleven, Landais, and Leite-Mariante (2024) extend this design to 134 countries and show that the size of the child penalty correlates strongly with local gender-norm measures. This situates any single-country estimate (including the Israeli one below) within a broader cross-national distribution rather than treating it as a universal constant. Goldin (2014) offers the leading explanation for \emph{why} the penalty persists even where formal barriers to women's employment have fallen: many high-paying jobs are ``greedy,'' rewarding long, inflexible, and unpredictable hours disproportionately. Workers who need temporal flexibility -- disproportionately mothers -- pay a wage penalty for it. This framework generates a sharp, testable prediction. An exogenous change that lowers the cost of supplying flexibility (working from home) should differentially help mothers, either by raising their employment probability, their hours, or both.

\paragraph{Does remote work narrow the gap in practice?} The evidence on whether flexible-work arrangements realize this potential is mixed. In a field experiment on Italian white-collar workers, Angelici and Profeta (2024) find that randomly assigned flexible (``smart'') working arrangements shifted the division of household tasks. Women's reported well-being also increased, suggestive of a genuine flexibility benefit. Real-time survey evidence from the acute pandemic period tells a more ambiguous story. Adams-Prassl et al. (2020) show that WFH-feasible jobs were substantially better shielded from the 2020 employment and earnings shock. Women in WFH-feasible jobs, however, simultaneously absorbed a larger share of the additional childcare burden created by school and daycare closures. The same shock that protected employment could, in principle, degrade the quality of that employment for mothers specifically. Alon et al. (2020) formalize the macro logic for why COVID-19-type shocks hit mothers asymmetrically (sector composition plus school closures). They do not take a stand on whether WFH exposure specifically offsets or compounds this asymmetry. This paper is designed to adjudicate empirically between these two directions using post-acute-period (2021--2023) data, rather than the initial-shock period that most of this literature studies.

\paragraph{Measuring WFH exposure, and the Israeli setting.} Since realized WFH status is itself a labor-market outcome -- jointly determined with employment and hours -- causal designs need a measure of WFH \emph{exposure} that is fixed before the outcome is realized. Dingel and Neiman (2020) provide the standard tool: an occupation-level teleworkability score built from U.S. task-content data. This score is not a perfect proxy outside the U.S. institutional context. As documented below, it substantially over-predicts realized WFH take-up in occupations where Israeli institutional practice (e.g., the Ministry of Education's in-person schooling policy) diverges from the U.S. task-content prediction. For the Israeli labor market specifically, Yaish, Mandel, and Kristal (2021) document that the 2020 lockdown reallocated paid and unpaid labor within dual-earner households in gender-unequal ways. Madhala and Bental (2021) provide a descriptive occupation-level baseline of WFH feasibility among Israeli workers. Buzaglo-Baris (2023) shows that within-firm sorting, not occupation choice alone, is a first-order driver of the Israeli gender wage gap. This is a reminder that any occupation-based exposure measure captures only part of the relevant heterogeneity.

\paragraph{Added value of this paper.} Relative to this literature, the present analysis makes four contributions. First, it uses population-scale administrative-style survey microdata spanning both pre- and multiple post-COVID years (2017--2019 and 2021--2023), rather than a single cross-sectional pandemic survey. This allows separate identification of extensive- and intensive-margin responses -- a distinction the WFH-and-gender literature above largely does not draw. Second, it instruments the mechanism with a \emph{pre-determined}, occupation-based exposure measure rather than realized WFH status. This avoids the reverse-causality problem inherent in treating self-reported WFH as an explanatory variable for employment or hours. Third, it exploits an Israeli institutional peculiarity: a furlough (\textit{Halat}) system that keeps the formal employment relationship intact during disruptions, unlike layoff-based adjustment in most U.S.-focused studies. As shown in Section~5, this has first-order consequences for \emph{which} margin absorbs any labor-supply response. Fourth, and most directly, it is among the first studies to test the WFH-exposure mechanism specifically in Israel's post-acute-pandemic labor market.

\section{Methodology}

\subsection{Hypotheses}

Two hypotheses follow directly from the flexibility framework in Goldin (2014) and the mixed empirical record above. \textbf{H1 (extensive margin):} occupational WFH exposure raises mothers' employment probability relative to other workers after 2021, if flexibility relaxes a binding participation constraint. \textbf{H2 (intensive margin):} even absent a detectable extensive-margin effect, WFH exposure raises mothers' hours conditional on employment, if the participation margin is structurally protected (e.g., by a furlough regime) while hours remain the flexible margin of adjustment.

\subsection{Empirical strategy}

The baseline specification is a difference-in-differences comparing mothers to other working-age women before and after the COVID-19 shock:
\begin{equation}
Employed_{it} = \beta_0 + \beta_1 Mother_i + \beta_2 Post_t + \beta_3 (Mother_i \times Post_t) + X_{it}'\gamma + \varepsilon_{it}
\label{eq:baseline}
\end{equation}
where $Post_t=1$ for survey years 2021--2023 and $0$ for 2017--2019 (2020 is excluded entirely; no raw survey extract exists for that year), and $X_{it}$ is the control vector defined in Section~4. The coefficient of interest, $\beta_3$, is estimated on the pooled sample and separately on Jewish and Arab subsamples, with standard errors clustered on the individual (\texttt{IDPUF}). An identical specification with $WorkHoursCont_{it}$ (usual weekly hours) as the dependent variable, restricted to $Employed_{it}=1$, identifies the intensive-margin analogue of $\beta_3$.

The mechanism test adds a pre-determined, occupation-based WFH-exposure regressor in a triple-difference (DDD) specification:
\begin{align}
Employed_{it} = {}& \beta_0 + \beta_1 Mother_i + \beta_2 Post_t + \beta_3 WFHExposure_i \notag \\
& + \beta_4 (Mother_i\times Post_t) + \beta_5(Mother_i \times WFHExposure_i) \notag \\
& + \beta_6(Post_t\times WFHExposure_i) + \beta_7 (Mother_i \times Post_t \times WFHExposure_i) \notag \\
& + Mother_i\times GilNK_i + X_{it}'\gamma + \varepsilon_{it}
\label{eq:ddd}
\end{align}
The coefficient of interest is $\beta_7$. $WFHExposure_i$ is constructed entirely from \emph{pre-period} (2017--2019) occupational composition (Section~4), so it varies at the level of a demographic cell rather than the individual. Standard errors are accordingly clustered at that cell level (the interaction of age group, education, and district) rather than at the individual level. This avoids understating the standard error on a regressor with no within-cluster variation (a Moulton-type correction). Olden and Møen (2022) formalize why this triple-difference design is attractive relative to running two separate difference-in-differences. It requires only a single parallel-trends assumption on the triple interaction, rather than two independent parallel-trends assumptions that would need to hold simultaneously.

The $Mother_i \times GilNK_i$ term (an interaction with the categorical age-group control) is included because a pre-period covariate balance check reveals a real age imbalance between mothers and non-mothers. This imbalance's \emph{magnitude} varies systematically with WFH-exposure quartile. The gap is $-0.797$ age-group units in the lowest exposure quartile ($t\approx-80.9$), narrowing and reversing to $+0.239$ in the highest quartile ($t\approx21.2$; $n=201{,}388$ pre-period observations) -- a pattern an additive age control cannot absorb.

\subsection{Identification assumptions and threats}

The design's central identifying assumption is parallel trends: absent the WFH-exposure shock, the employment (or hours) gap between mothers and comparison workers would have evolved similarly across exposure levels. This is tested with an event-study specification (Mother $\times$ year, 2019 omitted) and a joint Wald test on the pre-2020 Mother-by-year coefficients. Both are consistent with flat pre-trends in every specification reported below, including the child-age heterogeneity specification in Section~5.

Three further threats are addressed directly. First, because the intensive-margin regression conditions on $Employed_{it}=1$, and WFH exposure is hypothesized to affect employment itself, conditioning on employment risks selecting on a mediator. This is addressed with Lee (2009) trimming bounds, adapted to the difference-in-differences setting following the standard procedure. These bounds constrain the intensive-margin effect under a monotone-selection assumption rather than assuming selection away. Second, an individual-level panel fixed-effects design was considered and rejected: only 1.9\% of sampled individuals are re-interviewed across the pre/post boundary, far too few for within-person identification to be credible. Third, the exposure regressor's relevance is verified directly: pre-period $WFHExposure_i$ strongly and stably predicts realized WFH status in the post period. This holds across every iteration of the exposure-cell construction described in Section~4 (coefficients of 0.93 to 1.79, all $p<0.01$), ruling out the possibility that a null triple-difference simply reflects a non-informative regressor.

\section{Data}

The data are Israeli Central Bureau of Statistics Labor Force Survey microdata for survey years 2017--2019 and 2021--2023 (2020 is unavailable and excluded). The analysis sample is restricted to women aged 25--59; a parallel men's subsample is constructed for a gender-placebo check. The core regression sample contains $n=364{,}784$ observations (Jewish subsample $n=283{,}975$; Arab subsample $n=65{,}598$). The primary triple-difference sample, after the exposure-cell merge, contains $n=355{,}984$ observations clustered into roughly 210 age-education-district cells.

\paragraph{Outcome variables.} $Employed_i=1$ iff the CBS collapsed employment code $Muasak=1$. A data-quality issue required a correction here: $Muasak=1$ includes workers on furlough (\textit{Halat}, $SibaNeedar=9$) under standard ILO ``employed, temporarily absent'' conventions. The furlough rate rose sharply from a 0.1--0.3\% pre-period baseline to 3.27\% in 2021 and 1.20\% in 2023. All headline results below use the corrected outcome (furlough-excluded, reported as $Employed$); results under the uncorrected definition are noted where they differ materially. $WorkHoursCont_i$ is usual weekly hours, imputed from CBS's banded hours variable via within-period bin medians.

\paragraph{Treatment variables.} $Mother_i=1$ iff the respondent reports any child under 17. $Post_t=1$ for 2021--2023.

\paragraph{WFH-exposure measure.} $WFHExposure_i$ is a shift-share measure built entirely from pre-period (2017--2019) data. Each demographic cell's exposure score is the survey-weighted average of a \emph{calibrated} occupation-level teleworkability score, using that cell's own pre-period occupational composition. The occupation-level score starts from the Dingel and Neiman (2020) crosswalk. It is calibrated downward where the external U.S.-based score is statistically distinguishable from realized 2021--2023 Israeli WFH take-up in that occupation. For example, teaching has an external score of 0.966 versus realized Israeli usage of 0.063, reflecting the Ministry of Education's in-person schooling policy. Because this measure is defined for every respondent regardless of employment status, it avoids conditioning the design's third difference on employment itself. The exposure-cell partition -- age group $\times$ education $\times$ district $\times$ marital status $\times$ religiosity $\times$ birth continent, 8{,}884 cells -- is deliberately built on a \emph{finer} variable set than the regression's own controls. This avoids the loss of statistical power that results when the exposure regressor is aliased with the regression's own fixed effects (documented in Section~5).

\paragraph{Controls.} Every specification controls for marital status, religiosity, a five-category age group, district of residence, and a six-category education variable. Age enters only as this categorical group because no continuous age or birth-year variable exists anywhere in the raw CBS extract for any sample year.

\paragraph{Limitations.} CBS survey weights ($MishkalSofi$, etc.) are not applied in any outcome regression reported below; all coefficients are unweighted sample estimates and should not be read as population-representative. This is a deliberate, documented scope decision rather than an oversight, made to keep the identification strategy legible. A weighted robustness check exists in the underlying codebase but has not been run, pending a separate methodological decision. The Lee bounds correct only for excess selection into the $Mother=1,\,Post=1$ employment cell, not for possible under-selection in that cell.

\section{Findings}

The results are presented in four steps: baseline extensive- and intensive-margin effects of motherhood (Section~5.1), the primary test of whether WFH exposure explains either pattern (Section~5.2), a broader search for extensive-margin heterogeneity once that test proves inconclusive (Section~5.3), and a synthesis of the overall pattern (Section~5.4).

\subsection{Baseline effects: employment is flat, hours rise}

Table~\ref{tab:baseline} reports the baseline difference-in-differences on employment. The pooled coefficient on $Mother\times Post$ is small and statistically indistinguishable from zero ($-0.0053$, SE $0.0052$). This null holds separately in the Jewish ($-0.0041$, SE $0.0058$) and Arab ($-0.0208$, SE $0.0138$) subsamples. A direct test of equality between the two subsample estimates likewise fails to reject equality ($-0.0122$, SE $0.0155$, $t\approx-0.79$). In sum, no detectable change in mothers' employment probability relative to other working-age women after 2021 is observed, either in the full population or in either major population subgroup.

\begin{table}[htbp]
\centering
\caption{Baseline difference-in-differences: employment}
\label{tab:baseline}
\begin{tabular}{lccc}
\hline
 & Pooled & Jewish & Arab \\
\hline
Mother $\times$ Post & $-0.0053$ & $-0.0041$ & $-0.0208$ \\
 & $(0.0052)$ & $(0.0058)$ & $(0.0138)$ \\
Post & $0.0153^{***}$ & & \\
 & $(0.0043)$ & & \\
\hline
$N$ & 364{,}784 & 283{,}975 & 65{,}598 \\
\hline
\end{tabular}
\par\smallskip
{\footnotesize Dependent variable: $Employed$ (furlough-corrected). Standard errors clustered on individual, in parentheses. Controls: marital status, religiosity, age group, district, education. $^{*}p<0.10$, $^{**}p<0.05$, $^{***}p<0.01$.}
\end{table}

The intensive margin exhibits a markedly different pattern (Table~\ref{tab:intensive}). Conditional on employment, mothers' hours rise by $0.83$ hours/week relative to other workers after 2021 ($p<0.01$) -- roughly 2\% of a typical 40-hour week. This is a modest but economically meaningful shift, and among the most precisely estimated coefficients reported here. Because conditioning on employment risks selecting on a mediator if WFH pulls marginal mothers into work, Lee~(2009) trimming bounds are reported alongside the point estimate. The bounded range is $[0.213,\,1.313]$, entirely positive, though the lower bound alone is not itself statistically distinguishable from zero. The effect is therefore positive in direction, with some residual uncertainty about its precise lower-bound magnitude.

\begin{table}[htbp]
\centering
\caption{Intensive margin: hours worked, conditional on employment}
\label{tab:intensive}
\begin{tabular}{lc}
\hline
 & Mother $\times$ Post \\
\hline
Point estimate & $0.826^{***}$ \\
 & $(0.183)$ \\
Lee (2009) lower bound & $0.213$ \\
Lee (2009) upper bound & $1.313^{***}$ \\
\hline
$N$ & 281{,}750 \\
\hline
\end{tabular}
\par\smallskip
{\footnotesize Dependent variable: usual weekly hours, $Employed=1$ sample only. Standard errors clustered on individual, in parentheses. $^{***}p<0.01$.}
\end{table}

A gender-placebo check supports a motherhood-specific interpretation. The equivalent Father $\times$ Post interaction in the male subsample is small and statistically null under both DDD specifications ($-0.0016$, SE $0.0216$; $-0.0098$, SE $0.0222$; $n=331{,}788$). There is thus no symmetric pattern among fathers that would suggest the (also null) baseline estimate reflects something other than a motherhood-specific channel.

\subsection{The WFH-exposure mechanism: underpowered, not cleanly zero}

Table~\ref{tab:ddd} reports the primary triple-difference test linking employment to WFH exposure specifically. The coefficient of interest, $Mother\times Post\times WFHExposure$, is small and statistically null in both specifications ($-0.0257$, SE $0.0725$; $-0.0194$, SE $0.0722$). This null, however, should not be read as evidence against the flexibility mechanism. The design's minimum detectable effect, even after two successive refinements to the exposure-cell construction that cut it by roughly half, remains approximately 26\% of the baseline employment rate. This is a threshold well above what a plausible true effect would need to clear to be detected. The $Post\times WFHExposure$ main effect, by contrast, is robustly negative and statistically significant throughout every iteration of the design ($-0.15$, $p<0.05$). This indicates a general post-2021 employment decline in high-exposure occupations that is unrelated to motherhood specifically.

\begin{table}[htbp]
\centering
\caption{Primary triple-difference: employment and WFH exposure}
\label{tab:ddd}
\begin{tabular}{lcc}
\hline
 & Additive & Interacted cell FE \\
\hline
Mother $\times$ Post $\times$ WFHExposure & $-0.0257$ & $-0.0194$ \\
 & $(0.0725)$ & $(0.0722)$ \\
Post $\times$ WFHExposure & $-0.152^{**}$ & $-0.152^{**}$ \\
\hline
Minimum detectable effect & $0.203$ & $0.202$ \\
$N$ & 355{,}984 & 355{,}984 \\
\hline
\end{tabular}
\par\smallskip
{\footnotesize Dependent variable: $Employed$. Standard errors clustered on the (age group $\times$ education $\times$ district) exposure cell, in parentheses. Both specifications include $Mother\times GilNK$. $^{**}p<0.05$.}
\end{table}

Two further diagnostics merit brief discussion. First, adding the $Mother\times GilNK$ term surfaces a genuine, if secondary, finding unrelated to WFH exposure. The interacted-FE specification's base $Mother$ coefficient moves from an insignificant $-0.001$ to a significant $-0.077$ ($p<0.01$) once interacted with age group. The interaction terms for the three oldest age groups are all positive and significant, consistent with an age-widening motherhood employment gap that a simple additive age control obscures. Second, the intensive-margin analogue of the DDD is marginally positive ($+4.61$ hours, SE $2.64$, $p<0.10$). This suggests, without reaching conventional significance, that the robust intensive-margin effect documented above may be partly WFH-mediated even where the extensive-margin DDD fails to detect a corresponding effect.

The furlough correction discussed in Section~4 moves the primary DDD estimates \emph{closer} to zero, not away from it ($-0.0257\rightarrow-0.0150$ and $-0.0194\rightarrow-0.0079$). The underlying mechanism is orthogonal to motherhood: the furlough rate declines monotonically with WFH-exposure quartile for mothers and non-mothers alike. There is no differential pattern by motherhood status within quartile. The correction therefore cleans up a genuine measurement problem without altering the substantive conclusion on the mechanism.

\subsection{A systematic search for extensive-margin heterogeneity}

Given the flat pooled extensive-margin null, five further heterogeneity dimensions were tested, covering education, single parenthood, a combined hours-including-zeros margin, occupational sorting, and youngest child's age. Four return clean nulls: education (joint Wald $F(5,79069)=1.19$, $p=0.311$) and single parenthood ($F(2,79069)=0.61$, $p=0.543$). Occupational sorting into more- or less-exposed occupations among the already-employed is likewise null ($0.0006$, SE $0.0034$). The combined margin shows a marginal positive coefficient that decomposes almost exactly as the mechanical sum of the already-documented intensive-margin effect and the near-zero extensive-margin one, so it carries no independent information.

\begin{table}[htbp]
\centering
\caption{Employment by youngest child's age}
\label{tab:childage}
\begin{tabular}{lc}
\hline
Child age group $\times$ Post & Coefficient (SE) \\
\hline
0--1 & $-0.0146^{*}$ (0.0076) \\
2--4 & $-0.0155^{**}$ (0.0072) \\
5--17 & null \\
\hline
Joint Wald test & $F(5,79069)=2.02$, $p=0.072$ \\
$N$ & 364{,}784 \\
\hline
\end{tabular}
\par\smallskip
{\footnotesize Dependent variable: $Employed$. Reference category: no children under 18. $^{*}p<0.10$, $^{**}p<0.05$.}
\end{table}

The fifth dimension, youngest child's age, is the one robust extensive-margin finding of the paper (Table~\ref{tab:childage}). Mothers of children aged 0--1 and, more precisely, 2--4 show a significant relative decline in employment after 2021 ($-0.0146$, $p<0.10$; $-0.0155$, $p<0.05$), while mothers of older children (5--17) show no effect. This pattern survives three separate robustness checks. First, it is essentially unchanged when an explicit age-confound interaction is added ($-0.0146\rightarrow-0.0142$; $-0.0155\rightarrow-0.0153$). Second, it strengthens, rather than weakens, under the furlough-corrected outcome ($-0.0155\rightarrow-0.0169$, $p<0.05$). Third, its timing pattern is more consistent with a structural shift than a transient pandemic-year disruption: it persists and grows through 2021--2023 rather than peaking in 2021 and fading. Flat pre-trends are confirmed by the same event-study check described in Section~3.

Critically, this finding cannot be linked to WFH exposure: adding the triple interaction of child-age group, $Post$, and $WFHExposure$ returns a clean null (joint Wald $F(5,209)=0.78$, $p=0.565$, $n=355{,}984$). The employment decline among mothers of young children after 2021 is real and robust, but its mechanism -- WFH-related or otherwise -- remains an open question that this design cannot resolve.

\subsection{Synthesis}

Taken together, these results locate the post-2021 labor-supply adjustment for mothers almost entirely on the intensive margin: hours rise significantly and robustly. Employment itself remains essentially flat, except for a specific, age-targeted decline among mothers of very young children that is not attributable to WFH exposure by any test conducted here. One institutional account is consistent with this pattern: Israel's furlough (\textit{Halat}) system preserves the formal employment relationship through disruption at comparatively low cost to the worker. Labor-supply adjustment is thus likely to occur first, or exclusively, on the less costly margin -- hours -- before reaching the more costly margin of exiting employment altogether. A purely statistical account is not mutually exclusive: the extensive-margin DDD requires an implausibly large true effect (on the order of a quarter of the baseline employment rate) to be detectable at all. The continuous hours outcome, by contrast, carries strictly more information per observation and is correspondingly better powered to detect an effect of realistic magnitude.

\section{Conclusion}

% ponytail: per the assignment guidelines, this section must be written
% individually by the student -- this is a fully-written MODEL answer,
% requested explicitly as a benchmark to compare your own wording and
% analysis against, not a section meant to survive into the submission.
% Replace it with your own conclusion before submitting.

This paper asked whether occupational exposure to work-from-home opportunities altered mothers' position in the Israeli labor market after COVID-19. The evidence yields an asymmetric answer: on one margin, robustly so; on the margin the underlying mechanism was designed to explain, the data remain inconclusive. Mothers' employment probability, relative to other working-age women, is statistically unchanged after 2021 -- a null that holds in the pooled sample, within the Jewish and Arab subsamples separately, and under a gender placebo. Conditional on employment, by contrast, mothers' hours rise by nearly a full hour per week, a shift that is both statistically precise and economically non-trivial at national scale. The most plausible interpretation is that Israeli mothers did not gain a new margin of entry into or exit from employment after the pandemic, but that many already employed adjusted the intensity of their labor supply.

This pattern merits attention independent of whether WFH exposure proves to be the underlying mechanism. It does not conform to what a straightforward application of Goldin's (2014) framework would predict: if flexibility relaxes a binding participation constraint, the extensive margin might be expected to move first. Instead, the results are more consistent with a setting in which the participation margin is institutionally rigid: Israel's furlough (\textit{Halat}) regime allows the firm-worker relationship to survive a demand shock without formal separation, so that the marginal unit of labor supply subject to adjustment is an hour rather than a job. If this interpretation is correct, it carries a direct policy implication: expanding WFH availability should be evaluated as a lever for improving the quality and intensity of employment among mothers already in the labor force, rather than as an instrument expected to raise labor-force participation. The evidence here offers limited support for the view that WFH mandates would meaningfully raise Israel's female labor-force participation rate, but stronger support for the view that they would ease the burden on mothers already employed.

The one margin on which employment does move is also the one this design cannot fully explain. Mothers of children aged 0--4 exhibit a real, robust decline in employment after 2021 that survives every falsification check applied here -- an age-confounding check, the furlough correction, and a pre-trend and timing test -- yet is unrelated to occupational WFH exposure. This constitutes the paper's most consequential negative result: it rules out the most immediate causal account (that WFH exposure differentially reduced employment among mothers of young children) without identifying a replacement. Two candidate explanations appear most plausible and are, in principle, testable with data beyond the scope of the present study: a childcare-supply channel, in which disruptions to formal daycare and preschool availability outlasted the acute 2020 shock, and a composition or selection channel operating through which mothers of young children were already attached to the labor force before the pandemic. Distinguishing between these would require linking the CBS survey to childcare-availability data or to a source with direct measures of unmet childcare demand, a natural extension of this analysis.

Two design limitations bound the scope of these conclusions. First, the WFH-exposure mechanism test is underpowered rather than cleanly negative: even after two successive refinements to the exposure-cell construction, the design can rule out only effects on the order of a quarter of the baseline employment rate or larger, a demanding threshold. A denser panel, ideally one that follows the same individuals across the pre/post boundary at a rate well above the 1.9\% observed here, would both sharpen this test and reopen the individual fixed-effects design set aside in Section~3 for lack of a usable panel. Second, every estimate reported here is unweighted; a population-representative version of this analysis, incorporating CBS survey weights, constitutes a deliberately deferred but necessary robustness check before these findings should inform policy.

Considered jointly, the results argue for a more disaggregated view of how labor markets respond to flexibility shocks than a single question -- whether WFH helped mothers -- can capture: the answer depends on which margin, which age of child, and which institutional safeguards are already in place. This disaggregation, more than any individual point estimate reported above, constitutes the paper's principal contribution.

\newpage
% ponytail: no Hebrew-language sources are cited in this draft, so the
% "separate Hebrew section at the very start of the bibliography" rule has
% nothing to place yet -- add that section first if a Hebrew source is
% added later. This list is hand-typed (matching the in-text citations,
% also hand-typed -- see the note near \usepackage{hyperref} above) to
% match the assignment's specific name/ordering rules exactly, which no
% standard citation-package style produces automatically.
\section*{References}
\begin{list}{}{\leftmargin=1.5em \itemindent=-1.5em}

\item Adams-Prassl, Abi, Teodora Boneva, Marta Golin, and Christopher Rauh (2020). ``Inequality in the Impact of the Coronavirus Shock: Evidence from Real Time Surveys.'' \textit{Journal of Public Economics}, 189, 104245.

\item Alon, Titan, Matthias Doepke, Jane Olmstead-Rumsey, and Mich\`{e}le Tertilt (2020). ``The Impact of COVID-19 on Gender Equality.'' NBER Working Paper No.\ 26947.

\item Angelici, Marta, and Paola Profeta (2024). ``Smart Working: Work Flexibility Without Constraints.'' \textit{Management Science}, 70(3), 1680--1705.

\item Barrero, Jose Maria, Nicholas Bloom, and Steven J. Davis (2021). ``Why Working from Home Will Stick.'' NBER Working Paper No.\ 28731.

\item Buzaglo-Baris, Shira (2023). ``Firm Effect and the Israeli Gender Wage Gap.'' Bank of Israel Discussion Paper No.\ 2023.17.

\item Dingel, Jonathan I., and Brent Neiman (2020). ``How Many Jobs Can Be Done at Home?'' \textit{Journal of Public Economics}, 189, 104235.

\item Goldin, Claudia (2014). ``A Grand Gender Convergence: Its Last Chapter.'' \textit{American Economic Review}, 104(4), 1091--1119.

\item Kleven, Henrik, Camille Landais, and Jakob Egholt S\o{}gaard (2019). ``Children and Gender Inequality: Evidence from Denmark.'' \textit{American Economic Journal: Applied Economics}, 11(4), 181--209.

\item Kleven, Henrik, Camille Landais, and Gabriel Leite-Mariante (2024). ``The Child Penalty Atlas.'' \textit{Review of Economic Studies}, 92(5), 3174--3216.

\item Lee, David S. (2009). ``Training, Wages, and Sample Selection: Estimating Sharp Bounds on Treatment Effects.'' \textit{Review of Economic Studies}, 76(3), 1071--1102.

\item Madhala, Shavit, and Benjamin Bental (2021). ``The Ability to Work from Home Among Workers in Israel.'' Taub Center for Social Policy Studies in Israel.

\item Olden, Andreas, and Jarle M\o{}en (2022). ``The Triple Difference Estimator.'' \textit{The Econometrics Journal}, 25(3), 531--553.

\item Yaish, Meir, Hadas Mandel, and Tali Kristal (2021). ``Has the Economic Lockdown Following the Covid-19 Pandemic Changed the Gender Division of Labor in Israel?'' \textit{Gender \& Society}, 35(2), 256--270.

\end{list}

\end{document}
